\chapter{العزم الحركي الكمومي}\label{ch:b3:quantum-angular-momentum}

وجّه مقراب راديو إلى رقعة مظلمة من السماء ووالفه على
\qty{115}{GHz}: يظهر خط ساطع حادّ كالإبرة --- إنها جزيئات أحادي
أكسيد الكربون، على بعد عشرة كلفنات فوق الصفر المطلق، تنزل درجة
واحدة من سلّم من الحالات \emph{الدورانية}. فتقلّبها، كأيّ شيء يدور
في الميكانيك الكمومي، مكمّى مرتين: إذ لا يمكن لطويلة العزم الحركي
إلا أن تكون $\sqrt{l(l+1)}\,\hbar$، ولا لمركبته على أيّ محور مختار
إلا $m\hbar$. ويستنتج هذا الفصل ذلك التكميم المزدوج --- مرة
جبريًا، انطلاقًا من المبدّلات الموروثة من أقواس بواسون في
\cref{ch:b3:hamiltonian-mechanics}، ومرة ملموسة، بوصفها
\omterm{prop:b3:quantum-angular-momentum:harmonics}{التوافقيات الكروية} التي تشكّل كل مدار ذري. وسيعطينا الجبر أكثر
ممّا نطلب: إذ يبيح قيمًا نصف صحيحة لا تحققها أيّ حركة مدارية،
وهو شغور تملؤه الطبيعة بعد فصلين بالسبين. وفي الطريق نلقى
سلالم الجزيئات الدورانية --- أي خطوط الموجات الملّيمترية التي يزن
بها الفلكيون الغاز البارد في المجرّات --- والعزوم المغناطيسية
التي أظهرت بها العزوم الحركية نفسها منفصمة أول مرة.

\section{جبر الدوران}

\begin{definition}[مؤثرات العزم الحركي]\label{def:b3:quantum-angular-momentum:operators}
\index{العزم الحركي!الكمومي}
العزم الحركي المداري هو المؤثر $\hat{\vect L} =
\hat{\vect r}\wedge\hat{\vect p}$، وبالمركبات $\hat L_z = \hat
x\hat p_y - \hat y\hat p_x$ وما يتبعه بالتبديل الدوري. وانطلاقًا من $[\hat x, \hat p_x] =
\iu\hbar$:
\[
[\hat L_x, \hat L_y] = \iu\hbar\,\hat L_z
\quad\text{(وبالتبديل الدوري)} , \qquad
[\hat L^2, \hat L_z] = 0 :
\]
فالمركبات غير متوافقة بعضها مع بعض --- ولا حالة تكون فيها اثنتان
منها حادّتين --- لكن المربع الكلي متوافق مع أيّ واحدة
منها: فالزوج $(\hat L^2, \hat L_z)$ هو الاختيار القياسي
للوسوم. وهذه المبدّلات هي بالضبط $\iu\hbar$ مضروبًا في أقواس
بواسون المحسوبة في \cref{ex:b3:hamiltonian-mechanics:brackets}:
أي قاموس ديراك في العمل.
\end{definition}

\begin{theorem}[الطيف، من الجبر وحده]\label{thm:b3:quantum-angular-momentum:spectrum}
\index{العدد الكمومي المغناطيسي}
لتكن $\hat J_x, \hat J_y, \hat J_z$ \emph{أيّ} ثلاثة مؤثرات
هرميتية تحقق علاقات التبادل أعلاه. عندئذٍ تكون القيم الذاتية
المشتركة للزوج $(\hat J^2, \hat J_z)$ هي
\[
\hat J^2:\ j(j+1)\hbar^2 , \qquad
\hat J_z:\ m\hbar , \quad m = -j, -j+1, \dots, +j ,
\]
حيث $j$ عدد \emph{صحيح أو نصف صحيح} غير سالب: أي $2j + 1$
قيمة للعدد $m$ من أجل كل $j$. \omterm{def:b3:harmonic-oscillator:ladder}{ومؤثرا السلّم} $\hat J_\pm = \hat
J_x \pm \iu\hat J_y$ ينقلان $m$ بالمقدار $\pm1$ عند $j$ ثابت:
\[
\hat J_\pm\ket{j,m} = \hbar\sqrt{j(j+1) - m(m\pm1)}\,\ket{j,m\pm1} .
\]
\end{theorem}

\begin{proof}[برهان جزئي]
$[\hat J_z, \hat J_\pm] = \pm\hbar\hat J_\pm$: فالتأثير بالمؤثر $\hat
J_\pm$ ينقل \omterm{def:b3:quantum-formalism:observable}{القيمة الذاتية} للمؤثر $\hat J_z$ بالمقدار $\pm\hbar$، عند
$\hat J^2$ ثابت (وهو يتبادل مع كل ما يُبنى من المؤثرات $\hat
J_i$). والنظيم $\|\hat J_\pm\ket{j,m}\|^2 = \hbar^2[j(j+1) -
m(m\pm1)]$ (المحسوب من $\hat J_\mp\hat J_\pm = \hat J^2 - \hat
J_z^2 \mp \hbar\hat J_z$) يجب أن يبقى غير سالب: فوجب أن ينتهي
السلّم من الأعلى عند $m_{\max}$ ما بحيث $m_{\max}(m_{\max}+1) =
j(j+1)$، أي $m_{\max} = j$، ومن الأسفل عند $m_{\min} = -j$.
والصعود من $-j$ إلى $+j$ بخطوات واحدية يفرض أن يكون $2j$ عددًا
صحيحًا غير سالب. وبذلك يتبرر ترميز \omterm{def:b3:quantum-formalism:observable}{القيمة الذاتية} $j(j+1)\hbar^2$
بعد وقوعه.
\end{proof}

\begin{remark}[شغور في الفهرس]\label{rem:b3:quantum-angular-momentum:halfinteger}
يبيح الجبر $j = \tfrac12, \tfrac32, \dots$ --- أي سلالم ذات
عدد زوجي من الدرجات. والحركة المدارية، كما سنبيّن تاليًا، لا
تستعمل إلا الأعداد الصحيحة. لكن الطبيعة لا تضيّع شيئًا: فالإلكترون
نفسه يحمل $j = \tfrac12$، بلا مدار وراءه
(\cref{ch:b3:spin-two-level})؛ والجبر المستنتَج هنا، بلا تغيير،
سيدير مغناطيسية الذرات وسبينات النوى والكيوبت.
\end{remark}

\section{العزم الحركي المداري: التوافقيات الكروية}

\begin{proposition}[التوافقيات الكروية]\label{prop:b3:quantum-angular-momentum:harmonics}
\index{التوافقيات الكروية}\index{العدد الكمومي المداري}
بالإحداثيات الكروية يكون $\hat L_z = -\iu\hbar\,\partial_\varphi$،
وتكون الدوال الذاتية المشتركة للزوج $(\hat L^2, \hat L_z)$ على الكرة
هي \emph{\omterm{prop:b3:quantum-angular-momentum:harmonics}{التوافقيات الكروية}} $Y_l^m(\theta, \varphi)$:
\[
\hat L^2\,Y_l^m = l(l+1)\hbar^2\,Y_l^m , \qquad
\hat L_z\,Y_l^m = m\hbar\,Y_l^m ,
\]
مع $l = 0, 1, 2, \dots$ \emph{صحيحًا} --- إذ تفرض أحادية قيمة
$\eu^{\iu m\varphi}$ أن يكون $m$ صحيحًا، ومنه $l$ صحيحًا. وأولها،
إلى معامل تنظيم: $Y_0^0 = \text{ثابت}$ (وهو الشكل $s$،
أي كرة)؛ و $Y_1^0 \propto \cos\theta$ و $Y_1^{\pm1}
\propto \sin\theta\,\eu^{\pm\iu\varphi}$ (وهي الأشكال $p$، ذات فصّين)؛
و $Y_2^0 \propto 3\cos^2\theta - 1$ وشريكاتها (وهي العائلة $d$).
وأما الزوجية: فاتجاه $Y_l^m(-\vect r$ العكسي$) = (-1)^l\,Y_l^m$.
\end{proposition}
\admitted

\begin{omfigure}
\begin{tikzpicture}[font=\small, scale=0.94]
  % s orbital
  \begin{scope}
    \draw[->] (0,-1.7) -- (0,1.8);
    \node[anchor=south] at (0,1.8) {$z$};
    \draw[omDef, thick, fill=omDef!18] (0,0) circle (1.1);
    \node at (0,-2.3) {$l = 0$ ($s$): كرة};
  \end{scope}
  % p_z
  \begin{scope}[xshift=4.4cm]
    \draw[->] (0,-1.7) -- (0,1.8);
    \node[anchor=south] at (0,1.8) {$z$};
    \draw[omThm, thick, fill=omThm!18, domain=0:360, smooth, samples=100]
      plot ({1.45*abs(cos(\x))*sin(\x)}, {1.45*abs(cos(\x))*cos(\x)});
    \node at (0,-2.3) {$l = 1$، $m = 0$: $|\cos\theta|$};
  \end{scope}
  % p_+-1
  \begin{scope}[xshift=8.8cm]
    \draw[->] (0,-1.7) -- (0,1.8);
    \node[anchor=south] at (0,1.8) {$z$};
    \draw[omExo, thick, fill=omExo!18, domain=0:360, smooth, samples=100]
      plot ({1.45*abs(sin(\x))*sin(\x)}, {1.45*abs(sin(\x))*cos(\x)});
    \node at (0,-2.3) {$l = 1$، $m = \pm1$: $|\sin\theta|$};
  \end{scope}
  % d_z2
  \begin{scope}[xshift=13.2cm]
    \draw[->] (0,-1.7) -- (0,1.8);
    \node[anchor=south] at (0,1.8) {$z$};
    \draw[omProp, thick, fill=omProp!18, domain=0:360, smooth, samples=140]
      plot ({0.62*abs(3*cos(\x)*cos(\x)-1)*sin(\x)}, {0.62*abs(3*cos(\x)*cos(\x)-1)*cos(\x)});
    \node at (0,-2.3) {$l = 2$، $m = 0$};
  \end{scope}
\end{tikzpicture}

{\small مخططات قطبية للمقدار $|Y_l^m(\theta)|$ (مقطع في مستوٍ
يحوي المحور $z$؛ والشكل الكامل هو مجسم الدوران).
وهذه الهياكل الزاوية، المستقلة عن أيّ كمون، ستكسو
كل ذرة في \cref{ch:b3:hydrogen-atom}.}
\end{omfigure}

\begin{example}[الدوّار الصلب وسلّمه]\label{ex:b3:quantum-angular-momentum:rotor}
لجزيء ثنائي الذرة يتقلّب بعزم عطالة $I$ \omterm{def:b3:hamiltonian-mechanics:hamiltonian}{هاميلتونية} $\hat H
= \hat L^2/2I$: وطاقاتها
\[
E_J = \frac{\hbar^2}{2I}\,J(J+1) = B\,J(J+1) ,
\qquad J = 0, 1, 2, \dots
\]
وكل منها منحلّ $(2J+1)$ مرة. ويحقق امتصاص الفوتون $\Delta J =
\pm1$، فتكون تواترات الامتصاص $\nu_{J+1\leftarrow J} =
2B(J+1)/h$: أي مشط من الخطوط \emph{متساوية التباعد} --- وهو التوقيع
الذي يُعرف به الطيف الدوراني من أول نظرة. ومن أجل أحادي
أكسيد الكربون، $B/h = \qty{57.6}{GHz}$: فيقع الخط الأساسي عند
\qty{115}{GHz} (أي $\lambda = \qty{2.6}{mm}$)، وهو حصان العمل في
علم الفلك الراديوي الملّيمتري
(\cref{pb:b3:quantum-angular-momentum:1}).
\end{example}

\section{استكمال الكمونات المركزية}

\begin{theorem}[الفصل في أيّ كمون مركزي]\label{thm:b3:quantum-angular-momentum:radial}
\index{المعادلة القطرية}\index{الحاجز الطاردي}
من أجل $V(r)$، يمكن أخذ الحالات المستقرة في الصورة
\[
\psi(\vect r) = \frac{u(r)}{r}\,Y_l^m(\theta, \varphi) ,
\]
حيث تحلّ $u$ \omterm{prop:b3:schrodinger-three-dimensions:swave}{المعادلة القطرية} الأحادية البعد
\[
-\frac{\hbar^2}{2m}\,u'' + \Big[V(r) +
\frac{\hbar^2\,l(l+1)}{2mr^2}\Big]u = E\,u , \qquad u(0) = 0 .
\]
وتُحلّ المسألة الزاوية مرة واحدة إلى الأبد بالتوافقيات $Y_l^m$؛ وكل
$l$ يضيف \emph{\omterm{thm:b3:quantum-angular-momentum:radial}{الحاجز الطاردي}} الدافع
$\hbar^2l(l+1)/2mr^2$ --- وهو الصورة الكمومية للكمون الفعّال في
\cref{ex:b3:lagrangian-mechanics:central} --- ويكون كل
مستوٍ من أجل $l$ معطى منحلًّا $(2l+1)$ مرة، لأن لا قوة مركزية
تعبأ باتجاه المحور $z$. وحيلة الموجة $s$ في
\cref{prop:b3:schrodinger-three-dimensions:swave} كانت
سطر $l = 0$ من هذه المبرهنة.
\end{theorem}

\begin{proof}[برهان جزئي]
ينقسم اللابلاسي في الصورة $\Delta = \tfrac1r\partial_r^2(r\,\cdot) -
\hat L^2/\hbar^2r^2$ (وهو مُسلَّم به؛ ومفاده أن $\hat L^2$
هو الجزء الزاوي من $-\hbar^2\Delta$). نُدخل $\psi = (u/r)Y_l^m$
ونستعمل \omterm{def:b3:quantum-formalism:observable}{القيمة الذاتية} للمؤثر $\hat L^2$: فنجد المعادلة المذكورة.
ويتبع الانحلال لأن $\hat H$ يتبادل مع المؤثرات الثلاثة $\hat
L_i$، ومؤثرات سلّمها تنقل $m$ من دون تغيير الطاقة.
\end{proof}

\begin{omfigure}
\begin{tikzpicture}
  \begin{axis}[omaxis, axis x line=bottom, axis y line=left,
    width=10.6cm, height=6.0cm, xmin=0, xmax=8, ymin=-1.35, ymax=1.3,
    xlabel={$r$ (بوحدة $a_0$)}, ylabel={$U_{\text{فعّال}}$ (بوحدة $E_{\text{I}}$)},
    xlabel style={at={(axis description cs:0.5,-0.1)}, anchor=north},
    xtick={2,4,6}, ytick=\empty,
    legend style={font=\scriptsize, at={(0.97,0.95)}, anchor=north east,
      draw=none, fill=none}]
    \addplot[omDef, thick, domain=1.35:8, samples=120] {-2/x};
    \addlegendentry{$l = 0$: كولوم صرف}
    \addplot[omThm, thick, domain=0.42:8, samples=160] {-2/x + 1/(x*x)};
    \addlegendentry{$l = 1$}
    \addplot[omExo, thick, domain=0.85:8, samples=160] {-2/x + 3/(x*x)};
    \addlegendentry{$l = 2$}
    \draw[gray, dashed] (axis cs:0,0) -- (axis cs:8,0);
  \end{axis}
\end{tikzpicture}

{\small الكمونات القطرية الفعّالة في مسألة كولوم: إذ يعزل
\omterm{thm:b3:quantum-angular-momentum:radial}{الحاجز الطاردي} $\hbar^2l(l+1)/2mr^2$ المبدأَ من أجل
$l \ge 1$. فلا تلامس النواةَ إلا الحالات $s$ --- وله نتائج
تمتد من الأطياف الذرّية إلى الأسر الإلكتروني المشعّ.}
\end{omfigure}

\section{العزوم المغناطيسية: العزم الحركي مرئيًا}

\begin{proposition}[العزم المغناطيسي المداري وأثر زيمان]\label{prop:b3:quantum-angular-momentum:magnetic}
\index{مغنطون بور}\index{أثر زيمان}
جسيم شحنته $q$ وكتلته $m$ وعزمه الحركي المداري
$\hat{\vect L}$ يحمل العزم المغناطيسي
\[
\hat{\vect\mu} = \frac{q}{2m}\,\hat{\vect L} ;
\]
ومن أجل الإلكترون، الوحدة الطبيعية هي \emph{\omterm{prop:b3:quantum-angular-momentum:magnetic}{مغنطون بور}}
$\mu_{\text{B}} = e\hbar/2m_{\text{e}} = \qty{5.79e-5}{eV/T}$
(قارن \cref{exo:b3:schrodinger-three-dimensions:5}). وفي مجال
$B\vect e_z$ تزيح الطاقة $-\hat{\vect\mu}\cdot\vect B$ كلَّ
مستوٍ بالمقدار $+m\,\mu_{\text{B}}B$ (وشحنة الإلكترون سالبة): فينفصم مستوٍ
ذو $l$ معطى إلى مركباته $2l + 1$ --- وهو
\emph{\omterm{prop:b3:quantum-angular-momentum:magnetic}{أثر زيمان}}، أول عرض مباشر لتكميم
$m$، وسبب تسمية $m$ \omterm{thm:b3:quantum-angular-momentum:spectrum}{بالعدد الكمومي المغناطيسي}.
وعند $B = \qty{1}{T}$ يكون الانفصام $\qty{58}{\micro eV}$: أي صغير
بجانب الطاقات الضوئية، وسهل التمييز بوصفه انزياحًا في الخطوط
الطيفية --- ومقروءًا بالعكس، مقياسًا مغناطيسيًّا: فانفصام زيمان
لخطوط ضوء الشمس يرسم خرائط المجالات المغناطيسية للبقع الشمسية.
\end{proposition}

\begin{proof}[برهان جزئي]
كلاسيكيًّا تكون شحنة على مدار حلقةَ تيار: $\mu = IA =
(qv/2\pi r)(\pi r^2) = qL/2m$؛ وترث العبارةُ المؤثرية ذلك
(وهي تتبع من التقارن مع $\hat{\vect A}$ في
\cref{prop:b3:lagrangian-mechanics:charge}). أما انزياح الطاقة فهو
نظرية الاضطرابات في المرتبة الأولى، متوقَّعًا هنا ومبرَّرًا في
\cref{ch:b3:perturbation-theory}.
\end{proof}

\begin{omfigure}
\begin{tikzpicture}[font=\small, >=Stealth]
  \draw[->] (0,-3.0) -- (0,3.15);
  \node[anchor=south] at (0,3.15) {$z$};
  % sphere radius sqrt6
  \pgfmathsetmacro\R{2.449}
  \draw[gray, dashed] (0,0) circle (\R);
  \foreach \m/\c in {2/omDef, 1/omThm, 0/omExo, -1/omThm, -2/omDef} {
    \pgfmathsetmacro\zz{\m}
    \pgfmathsetmacro\rr{sqrt(\R*\R-\m*\m)}
    \draw[\c, thick, ->] (0,0) -- (\rr,\zz);
    \draw[gray!60, thin, dashed] (-0.05,\zz) -- (\rr,\zz);
    \node[anchor=west, \c] at ({\rr+0.1},\zz) {$m = \m$};
  }
  \node[anchor=east, font=\footnotesize] at (-0.15,2.0) {$L_z = m\hbar$};
  \node[anchor=north, font=\footnotesize] at (1.3,-2.75) {$\|\vect L\| = \sqrt{6}\,\hbar$};
\end{tikzpicture}

{\small ``التكميم الفضائي'' من أجل $l = 2$: لمتجهة العزم الحركي
طول $\sqrt{l(l+1)}\,\hbar = \sqrt6\,\hbar$ لكنها لا تستطيع أن
تعرض على المحور $z$ غير الإسقاطات الخمسة $m\hbar$ --- ولا تعرض
طولها الكامل أبدًا: فلمّا كانت المركبات غير متوافقة، لا تستطيع
$\vect L$ أن تقع بالضبط على أيّ محور.}
\end{omfigure}

\begin{method}[العزم الحركي عمليًّا]\label{met:b3:quantum-angular-momentum:method}
(1) سِم الحالات بالزوج $(l, m)$ --- أو $(j, m)$ من أجل الجبر
المجرد --- ولا تطلب مركبتين معًا أبدًا. (2) والعناصر
المصفوفية: استعمل صيغ السلّم؛ و $\hat L_x = (\hat L_+ + \hat
L_-)/2$. (3) وفي كمون مركزي: اذكر $Y_l^m$، ولا تحلّ إلا
\omterm{prop:b3:schrodinger-three-dimensions:swave}{المعادلة القطرية} \omterm{thm:b3:quantum-angular-momentum:radial}{بالحاجز الطاردي}. (4) والطاقات الدورانية:
$BJ(J+1)$، والخطوط عند $2B(J+1)$ --- فاستخرج $B$، ومنه أطوال
الروابط، من التباعدات المتساوية. (5) والمسائل المغناطيسية: حوّل
العزوم الحركية إلى عزوم بمعدل $\mu_{\text{B}}$ لكل $\hbar$، والطاقات
بمعدل $\mu_{\text{B}}B$ لكل وحدة من $m$.
\end{method}

\section{تمارين}

\begin{exercise}[$\star$]\label{exo:b3:quantum-angular-momentum:1}
(أ) استنتج $[\hat L_x, \hat L_y] = \iu\hbar\hat L_z$ انطلاقًا من
المبدّلات القانونية. (ب) بيّن أن $[\hat L^2, \hat L_z] = 0$. (ج) بيّن أن
$[\hat L_z, \hat x] = \iu\hbar\hat y$: فماذا يولّد $\hat L_z$؟
(د) لماذا لا تقبل المركبات الثلاث أساسًا ذاتيًا مشتركًا
--- إلا من أجل حالة واحدة بعينها (وما هي)؟
\end{exercise}

\begin{exercise}[$\star$]\label{exo:b3:quantum-angular-momentum:2}
من أجل $l = 1$، في الأساس $\{\ket{1,1}, \ket{1,0}, \ket{1,-1}\}$:
(أ) اكتب مصفوفة $\hat L_z$؛ (ب) استعمل صيغة السلّم لتكتب
$\hat L_+$ و $\hat L_-$؛ (ج) ركّب $\hat L_x$ وتحقق من أن
قيمه الذاتية $\hbar, 0, -\hbar$؛ (د) حالة فيها $L_x =
+\hbar$ تُقاس على المحور $z$: أعطِ احتمالات النتائج
الثلاث.
\end{exercise}

\begin{exercise}[$\star$]\label{exo:b3:quantum-angular-momentum:3}
أحادي أكسيد الكربون: طول الرابطة \qty{0.113}{nm}، والكتلة المرجعة
$\qty{1.14e-26}{kg}$. (أ) احسب $I$ و $B = \hbar^2/2I$ بالميلي إلكترون-فولط.
(ب) احسب تواتر الخط $J = 1 \leftarrow 0$ وطول موجته.
(ج) احسب الخطين التاليين. (د) تقيس فلكية تباعد المشط
بخمسة أرقام: فأيّ مقدار جزيئي تحصل عليه، وبأيّ دقة؟
\end{exercise}

\begin{exercise}[$\star$]\label{exo:b3:quantum-angular-momentum:4}
جسيم في حالة $l = 2$. (أ) اذكر النتائج الممكنة
لقياس $L_z$. (ب) احسب أصغر زاوية بين $\vect L$ و
المحور $z$، باستعمال $\cos\theta = m/\sqrt{l(l+1)}$. (ج) لماذا
لا يمكن أن تكون الزاوية معدومة؟ (د) بيّن أن الزاوية الدنيا تؤول إلى
الصفر لمّا $l \to \infty$: فنستعيد المتجهات الكلاسيكية.
\end{exercise}

\begin{exercise}[$\star\star$]\label{exo:b3:quantum-angular-momentum:5}
الجزء الزاوي من $-\hbar^2\Delta$ هو $\hat L^2$، حيث
\[
\hat L^2 = -\hbar^2\Big[\frac{1}{\sin\theta}\,\partial_\theta
(\sin\theta\,\partial_\theta) +
\frac{1}{\sin^2\theta}\,\partial_\varphi^2\Big] .
\]
(أ) تحقق من أن $Y_1^0 \propto \cos\theta$ لها \omterm{def:b3:quantum-formalism:observable}{قيمة ذاتية} للمؤثر $\hat L^2$
مقدارها $2\hbar^2$. (ب) تحقق كذلك من $Y_1^{\pm1} \propto \sin\theta\,
\eu^{\pm\iu\varphi}$، ومن قيمها الذاتية للمؤثر $\hat L_z$. (ج)
تحقق من الزوجيات. (د) لماذا يجب أن يكون $\braket{Y_1^0}{Y_1^1} = 0$ من دون
حساب أيّ تكامل؟
\end{exercise}

\begin{exercise}[$\star\star$]\label{exo:b3:quantum-angular-momentum:6}
أيّ خط دوراني هو الأسطع؟ جمهرة المستوي $J$ هي
$\propto (2J+1)\,\eu^{-BJ(J+1)/k_{\text{B}}T}$. (أ) اشرح العاملين.
(ب) بيّن أن الأعظمية تقع قرب $J_{\max} \approx
\sqrt{k_{\text{B}}T/2B} - \tfrac12$. (ج) من أجل أحادي أكسيد الكربون عند \qty{10}{K}
(أي سحابة مظلمة) وعند \qty{300}{K}: أيّ الخطوط تسود؟ (د)
وبالقلب: سلّم أحادي أكسيد الكربون مشاهَدًا وذروته عند $J = 7$ يفضح أيّ
درجة حرارة؟
\end{exercise}

\begin{exercise}[$\star\star$]\label{exo:b3:quantum-angular-momentum:7}
\omterm{prop:b3:quantum-angular-momentum:magnetic}{أثر زيمان} العادي. خط طيفي عند \qty{500}{nm} يأتي من
انتقال $l = 2 \to l = 1$ في مجال $B = \qty{2}{T}$؛ وأهمل
السبين (وهو صالح من أجل حالات ``أحادية'' خاصة). (أ) ارسم المستويات
المنفصمة تقريبيًا. (ب) مع قاعدة الانتقاء $\Delta m = 0, \pm1$، بيّن أنه لا تظهر
غير \emph{ثلاثة} مواضع للخطوط، عند $0, \pm\mu_{\text{B}}B/h$.
(ج) احسب الانفصام بالجيغاهرتز وبالبيكومترات من طول الموجة.
(د) تنفصم معظم الخطوط الحقيقية إلى أكثر من ثلاث مركبات
(وهو زيمان ``الشاذّ''): فأيّ مكوّن ناقص، يحمله
الإلكترون نفسه، كان ذلك دليلًا تاريخيًا عليه؟
\end{exercise}

\begin{exercise}[$\star\star$]\label{exo:b3:quantum-angular-momentum:8}
الجدار الطاردي. من أجل الهدروجين ($V = -e^2/4\pi\varepsilon_0r$)،
اكتب $U_{\text{فعّال}}$ من أجل $l = 1$ بوحدتَي $E_{\text{I}}$ و
$a_0$. (أ) عيّن أصغريته وقيمتها. (ب) بيّن أن $U_{\text{فعّال}} >
0$ من أجل $r < \dots$ --- وأوجد نصف القطر الذي يسود الحاجز
داخله. (ج) قارن $|\psi|^2$ بجوار $r = 0$ في الحالات $s$ والحالات $p$
(مع $u \sim r^{l+1}$، وهو مُسلَّم به): فأيّ الحالات ``تلامس''
النواة؟ (د) الأسر الإلكتروني --- أي ابتلاع نواة أحدَ
إلكترونات ذرتها --- يجري في الغالب الأعمّ من الأغلفة $s$:
اشرح في جملة واحدة.
\end{exercise}

\begin{exercise}[$\star\star$]\label{exo:b3:quantum-angular-momentum:9}
على الحالة $\ket{l, m}$: (أ) بيّن أن $\langle\hat L_x\rangle = \langle\hat
L_y\rangle = 0$؛ (ب) بيّن أن $\langle\hat L_x^2\rangle = \langle\hat
L_y^2\rangle = \tfrac12[l(l+1) - m^2]\hbar^2$؛ (ج) تحقق من
\omterm{thm:b3:quantum-formalism:uncertainty}{علاقة الارتياب} $\Delta L_x\,\Delta L_y \ge \tfrac\hbar2
|\langle\hat L_z\rangle|$ على الحالة الممدودة $m = l$؛ (د)
فسّر صورة ``المخروط'' في الشكل على ضوء (أ) و
(ب).
\end{exercise}

\begin{exercise}[$\star\star\star$]\label{exo:b3:quantum-angular-momentum:10}
نطاقات الاهتزاز والدوران. انتقال اهتزازي تحت أحمر
($\hbar\omega_0$) في جزيء ثنائي الذرة يغيّر $J$ بالمقدار $\pm1$
في آن واحد. (أ) بيّن أن خطوط الامتصاص تقع عند $h\nu =
\hbar\omega_0 + 2B(J+1)$ (وهو الفرع $R$، من $J \to J+1$) وعند
$h\nu = \hbar\omega_0 - 2BJ$ (وهو الفرع $P$، من $J \to J-1$) مع $J \ge
1$. (ب) لماذا توجد \emph{فجوة} عند $\hbar\omega_0$ نفسه؟ (ج)
ارسم النطاق تقريبيًا: مشطان يحفّان بخط مركزي مفقود، وقوة كل
خط تتبع
\cref{exo:b3:quantum-angular-momentum:6}. (د) لنطاق
\qty{15}{\micro m} في \cref{pb:b3:harmonic-oscillator:1} هذا
التشريح بالضبط: فما الذي يحدّد عرضه الكلي، ومن ثمّ
``الجناحين'' اللذين يجعلان ثاني أكسيد الكربون المضاف فعّالًا؟
\end{exercise}

\begin{exercise}[$\star\star\star$]\label{exo:b3:quantum-angular-momentum:11}
إتمام الجبر. (أ) انطلاقًا من $\hat J_\mp\hat J_\pm = \hat J^2 -
\hat J_z^2 \mp \hbar\hat J_z$، احسب $\|\hat J_\pm\ket{j,m}\|^2$
ومعاملات السلّم. (ب) بيّن أنه لو لم ينتهِ السلّم
لصار نظيم ما سالبًا: واعرض الحالة المخالفة.
(ج) استنتج أن $m_{\max} - m_{\min} = 2j$ يجب أن يكون عددًا صحيحًا
غير سالب وأحصِ القيم المسموحة للعدد $j$. (د) أين بالضبط \emph{تعجز}
الحجة عن استبعاد أنصاف الأعداد الصحيحة --- وأيّ
شرط إضافي (أحادية القيمة على الكرة) يستبعدها
في الحركة المدارية وحدها؟
\end{exercise}

\begin{exercise}[$\star\star\star$]\label{exo:b3:quantum-angular-momentum:12}
أينشتاين--دي هاس: العزم الحركي ميكانيكي. أسطوانة حديدية
(نصف قطرها $a = \qty{1.0}{cm}$، وكثافتها $\rho = \qty{7.9e3}{kg/m^3}$، و
$n = \qty{8.5e28}{atoms/m^3}$) معلَّقة بليف فتل؛ وقلب
مغنطتها يقلب نحو $\hbar$ واحدًا من العزم الحركي لكل
ذرة. (أ) بالانحفاظ، يجب أن ترتدّ الأسطوانة: احسب العزم
الحركي المنقول في وحدة الحجم. (ب) مع عزم عطالة القضيب $\tfrac12 M
a^2$، بيّن أن القضيب يكتسب $\omega = 2n\hbar/\rho a^2$ وقدّره.
(ج) قادت تجربة 1915 القلبَ عند رنين الليف لتراكم
الركلات الضئيلة: قدّر سعة الاهتزاز
بعد $Q \sim 100$ قلبة رنينية، بأخذ قيمة
$\omega$ لكل ركلة. (د) خرجت النسبة المقيسة بين العزم المغناطيسي والعزم الحركي
ضعفَ التنبؤ المداري $e/2m_{\text{e}}$: فبمَ كان
ذلك العامل اثنان يبشّر؟
\end{exercise}

\section{مسألة: وزن المجرّات عند 2.6 ملّيمتر}

\begin{problem}[{مسألة نهاية الأسبوع --- أحادي أكسيد الكربون وعلم
الفلك الراديوي والكون البارد}]\label{pb:b3:quantum-angular-momentum:1}
معظم المادة التي تصنع النجوم في مجرّة هدروجين جزيئي بارد ---
وهو غير مرئي: فالجزيء H$_2$، وهو متناظر، لا ثنائي قطب له ولا
خطوط دورانية عند درجات حرارة السحب الباردة. وحلّ علم الفلك هو
رفيقه الوفيّ. فجزيء واحد من أحادي أكسيد الكربون لكل $\num{e4}$ من H$_2$، بسلّمه
عند \qty{115}{GHz}، يضيء كل سحابة جزيئية في
السماء. المعطيات: من أجل أحادي أكسيد الكربون، $B/h = \qty{57.6}{GHz}$ (أي $B =
\qty{0.238}{meV}$)؛ وطول الرابطة $r_0 = \qty{0.113}{nm}$؛ والكتلة
المرجعة $\mu = \qty{1.14e-26}{kg}$؛ و $k_{\text{B}}T$ عند \qty{10}{K} يساوي
\qty{0.862}{meV}.

\medskip\noindent\textbf{الجزء الأول --- الدوّار الكمومي.}
\begin{enumerate}
  \item انطلاقًا من $\hat H = \hat L^2/2I$، أعطِ الطاقات $E_J$ و
        انحلالاتها.
  \item تحقق من $B = \hbar^2/2I$ لأحادي أكسيد الكربون انطلاقًا من $r_0$ و $\mu$.
  \item مع قاعدة الانتقاء $\Delta J = \pm1$، استنتج
        تواترات الامتصاص $\nu_{J+1\leftarrow J} = 2B(J+1)/h$
        واذكر الثلاثة الأولى.
  \item لماذا تكون الخطوط \emph{متساوية التباعد} --- وماذا يعطي
        التباعد المقيس للفلكي (عبر $I$)؟
  \item احسب طول موجة الخط الأساسي، واشرح
        لماذا يحتاج رصده إلى مواقع مرتفعة جافة (فما الذي يمتصّ
        الأمواج الملّيمترية في غلافنا الجوّي --- تذكّر
        جارة جزيء \cref{pb:b3:harmonic-oscillator:1}،
        أي الماء)؟
  \item يحمل فوتون الخط $1 \to 0$ عزمًا حركيًا مقداره $\hbar$:
        فلماذا ليست $\Delta J = \pm1$ قاعدة تقريبية
        بل قانون انحفاظ؟
\end{enumerate}

\medskip\noindent\textbf{الجزء الثاني --- لماذا أحادي أكسيد الكربون لا H$_2$.}
\begin{enumerate}[resume]
  \item الجزيء H$_2$ متجانس النوى: اذكر لماذا لا يصدر إشعاعًا
        دورانيًا ثنائي القطب البتة.
  \item والجزيء H$_2$ \emph{خفيف} أيضًا: احسب الثابت
        الدوراني للجزيء H$_2$ ($\mu = m_{\text{p}}/2$ و $r_0 = \qty{0.074}{nm}$)
        وطاقة أول انتقال متاح له؛ وقارنها مع
        $k_{\text{B}}T$ عند \qty{10}{K}: فحتى بأبواب رباعي القطب
        الخلفية، هل كان H$_2$ البارد ليشعّ؟
  \item تكلّف أول درجة في أحادي أكسيد الكربون $2B = \qty{0.48}{meV}$ في مقابل
        $k_{\text{B}}T = \qty{0.86}{meV}$: فهل السلّم حيّ عند
        \qty{10}{K}؟
  \item احسب نسبة الجمهرتين $n_1/n_0$ عند \qty{10}{K}
        (مع إدخال الانحلال).
  \item قرب أيّ $J$ تبلغ الجمهرة ذروتها عند \qty{10}{K}
        (استعمل $J_{\max} \approx \sqrt{k_{\text{B}}T/2B} - \tfrac12$)؟
  \item لخّص في جملة واحدة عقد هذه المقايضة: ما الذي يقدّمه
        أحادي أكسيد الكربون، وما الذي يجب افتراضه في نسبته إلى
        H$_2$.
\end{enumerate}

\medskip\noindent\textbf{الجزء الثالث --- قراءة السماء.}
\begin{enumerate}[resume]
  \item يصل خط سحابة عند \qty{115.156}{GHz} بدل قيمة
        المخبر \qty{115.271}{GHz}: احسب سرعة السحابة
        على خط النظر، واتجاهها.
  \item الخط \emph{معرَّض} بأثر دوبلر بحركات داخلية مقدارها
        $\pm\qty{2}{km/s}$: احسب عرض الخط بالميغاهرتز، و
        قارنه بالعرض الحراري المنتظر عند \qty{10}{K}
        (حيث $v_{\text{th}} \sim \sqrt{k_{\text{B}}T/m_{\text{CO}}}
        \approx \qty{55}{m/s}$): فما الذي يسود التعريض؟
  \item يعطي رسم انزياح دوبلر لأحادي أكسيد الكربون عبر قرص مجرّة حلزونية
        منحني دورانها $v(r)$. وتبقى المنحنيات
        \emph{مسطحة} إلى ما وراء القرص المرئي بكثير: تذكّر (كتاب السنة
        الأولى، الثقالة) ماذا ينبغي أن يفعل $v(r)$ حول كتلة
        متمركزة، واذكر ما تقتضيه هذه الاستواء.
  \item نسبة سطوع الخط $2\to1$ إلى سطوع الخط
        $1\to0$ تقيس جمهرات المستويات: فأيّ مقدار فيزيائي
        واحد للسحابة تقدّمه؟
  \item يميّز مرصد ألما أحادي أكسيد الكربون في مجرّات غادر ضوءها والكون
        فتيّ؛ والتواتر المشاهَد للخط
        $1 \to 0$ من مجرّة عند ``انزياح أحمر $z = 1$''
        يصل عند نصف قيمته السكونية --- فعند أيّ تواتر، و
        (بالعودة إلى \cref{ex:b3:relativistic-kinematics:redshift})
        ما الذي مدّده؟
  \item انطلاقًا من إضاءة خط أحادي أكسيد الكربون المقيسة يذكر الفلكيون كتل
        السحب بالكتل الشمسية: اذكر سلسلة التحويلات (من فوتونات الخط
        $\to$ عدد جزيئات أحادي أكسيد الكربون $\to$ كتلة H$_2$) وأضعف
        حلقة فيها.
\end{enumerate}

\medskip\noindent\textbf{الجزء الرابع --- الحضّانات الباردة.}
\begin{enumerate}[resume]
  \item تقع أنوية تشكّل النجوم عند \qty{10}{K}: فبقانون فين
        يبلغ توهجها الحراري ذروته قرب \qty{290}{\micro m} ---
        وهو غير مرئي لأيّ مقراب ضوئي. في جملة واحدة: لماذا
        يكون السلّم الدوراني \emph{الميزان} الوحيد للحرارة
        وللمقياس الذي تقدّمه نواة كهذه؟
  \item يقع المستوي $J = 1$ فوق الحالة الأساسية بمقدار \qty{5.5}{K}
        (بوحدات الحرارة): فلماذا يجعل ذلك الخط $1\to0$
        مقياس حرارة بديعًا في نظام \qty{10}{K} بالضبط
        (بدل خط ضوئي مثلًا عند
        \qty{2}{eV} $\sim \qty{23000}{K}$)؟
  \item لماذا تكفّ نواة \emph{منهارة} في النهاية عن أن تكون
        مرئية في خط أحادي أكسيد الكربون ($1\to0$) --- انظر في ما تفعله الكثافة
        العالية والغبار بالخط وبإفلاته.
  \item تُظهر السحب الجزيئية أيضًا خطوط NH$_3$ و HCN وعشرات
        غيرها، لكل منها $B$ وثنائي قطب خاصّان: فما الذي يتيحه
        ثراء هذا ``القرص الراديوي الكيميائي'' للفلكيين أن يفكّكوه؟
  \item تتوهج السماء كلها عند \qty{2.7}{K} (وهي خلفية الميكروويف
        الكونية): فماذا يحلّ بتباين خطوط أحادي أكسيد الكربون في سحابة
        كلما اقتربت درجة حرارتها من
        \qty{2.7}{K}، ولماذا تكون تلك أرضية صلبة لهذا
        الميزان؟
  \item يجب أن يحفظ طبق ملّيمتري شكله في حدود
        $\lambda/20$: احسب ذلك التسامح عند \qty{115}{GHz}،
        وقارنه بالتسامح الذي يجب أن تحققه مرآة ضوئية
        (\qty{500}{nm}) --- ولماذا أمكن بناء أطباق ألما
        ذات الاثني عشر مترًا في الهواء الطلق؟
  \item لخّص النتيجة الموسومة: سلّم $BJ(J+1)$ مع $B/h =
        \qty{57.6}{GHz}$، معمورًا عند عشرة كلفنات ومقروءًا عند
        \qty{2.6}{mm}، هو كيف تُقاس كتلة الكون البارد ودرجة
        حرارته وسرعته ودورانه --- والدليل على المادة
        المظلمة حول المجرّات الحلزونية.
\end{enumerate}
\end{problem}
